\section{Implementation and Reproducibility}
\label{sec:impl}

\paragraph{Code organization}
The experiment pipeline is implemented under the project source directory with
modular components for configuration, data loading, and autoregressive baselines,
speculative decoding, metric aggregation, and runtime orchestration.
The notebook driver executes these modules in fixed phase order and writes
artifacts to the results directory.

\paragraph{Model loading and precision}
In the current configuration, target and draft models are loaded in 16-bit
floating point (FP16) for
the main sweep, with quantization controls retained in configuration for
alternative runs. The tokenizer and generation limits are kept constant between
baselines and speculative runs to preserve paired comparability.

\paragraph{Baselines}
We run one autoregressive baseline per regime (deterministic and stochastic)
with the same manifests and output-length constraints as speculative runs.
Each baseline artifact stores per-sample latency, output length, and task
predictions.

\paragraph{Speculative decoding implementation.}
The speculative loop follows standard token-level draft--verify logic.
At each cycle, the draft proposes up to $k$ tokens autoregressively; the
target verifies the proposal in one pass over the extended prefix; the accepted
prefix is committed up to the first mismatch; and one correction token is added
from the target distribution when needed. This procedure preserves target-level
generation semantics while exposing acceptance statistics needed for analysis.
Both draft and target states are tracked with key-value (KV) caches throughout
the decode loop.
Figure~\ref{fig:specdec-impl-flow} summarizes this implementation flow.

\paragraph{CUDA-graph runtime note}
In the Qwen3 run families, graph-capture metadata differs by hardware.
The RTX4090 fixed-policy matrix reports stable captured execution, whereas
RTX5090-laptop sensitivity runs show that graph-enabled toggles do not
consistently improve end-to-end latency. We therefore treat CUDA-graph usage
as an implementation detail rather than a standalone performance claim, and
base conclusions on end-to-end fixed-policy speedup trends.

\paragraph{Dynamic-scheduling bottleneck in native PyTorch paths}
Fallback eager execution should not be viewed only as an implementation miss.
It reflects a broader systems challenge for end-to-end dynamic speculation.
When acceptance-dependent control flow, sampling branches, and cache updates
change per step, preserving lossless semantics often requires CPU-visible
orchestration between kernels. Without a deeper runtime/kernel redesign, this
serial control overhead can erase expected gains from speculative parallelism,
as also emphasized in modern inference runtime roadmaps for dynamic
scheduling.

\paragraph{Tokenizer isomorphism as an enabling design choice}
Same-family draft/target pairing is a central engineering decision in this
study. Because the draft and target tokenizers are isomorphic, speculative
verification avoids per-step vocabulary alignment and token translation.
Cross-tokenizer pipelines introduce additional CPU-side mapping overhead that
can materially reduce or negate speculative speedup on constrained hardware.

\paragraph{Adaptive-$k$ implementation note}
For adaptive-$k$ runs, we reuse the same draft--verify core and switch only the policy
controller: $k$ is adapted online from recent acceptance behavior with guarded
fallback to stable paths when acceptance collapses. We treat the adaptive-$k$
run set as a runtime
policy layer over the fixed-policy implementation, not as a separate decoding
algorithm.

\begin{figure}[t]
\centering
\resizebox{0.98\columnwidth}{!}{%
\begin{tikzpicture}[
	node distance=4.8mm and 6mm,
	box/.style={rectangle, draw, rounded corners=1mm, align=center, minimum height=6mm, text width=31mm, inner sep=1.5mm},
	decision/.style={diamond, draw, align=center, aspect=2, inner sep=1.1mm, text width=22mm},
	line/.style={-Latex, thick, shorten >=1pt, shorten <=1pt}
]
\node[box] (start) {Start sample\\encode prompt and init KV cache};
\node[box, below=of start] (draft) {Draft model proposes\\up to $k$ tokens};
\node[box, below=of draft] (verify) {Target verifies draft block\\in one forward pass};
\node[decision, below=of verify] (accept) {All $k$ tokens\\accepted?};
\node[box, left=of accept, xshift=-9mm] (bonus) {Commit $k$ tokens\\+ one target bonus token};
\node[box, right=of accept, xshift=9mm] (correct) {Commit accepted prefix\\+ one corrected target token};
\node[box, below=of accept, yshift=-2mm] (crop) {Crop/update draft and target\\KV caches to committed prefix};
\node[decision, below=of crop] (stop) {EOS or max tokens\\reached?};
\node[box, below=of stop] (end) {Finish sample and\\record runtime/quality fields};

\draw[line] (start) -- (draft);
\draw[line] (draft) -- (verify);
\draw[line] (verify) -- (accept);
\draw[line] (accept) -- node[above, font=\scriptsize] {yes} (bonus);
\draw[line] (accept) -- node[above, font=\scriptsize] {no} (correct);
\draw[line] (bonus) |- (crop.west);
\draw[line] (correct) |- (crop.east);
\draw[line] (crop) -- (stop);
\draw[line] (stop) -- node[right, font=\scriptsize] {yes} (end);
\draw[line] (stop.west) -- ++(-14mm,0) |- node[pos=0.25, left, font=\scriptsize] {no} (draft.west);
\end{tikzpicture}%
}
\caption{Vanilla speculative decoding implementation flow.
The target model remains the verifier of record; accepted draft prefixes are
committed, and caches are cropped to preserve exact target-consistent decoding.}
\label{fig:specdec-impl-flow}
\end{figure}

\paragraph{Regimes and outputs}
The same verification core is used for deterministic and stochastic modes,
with only the sampling policy changed by regime parameters.
Each configuration writes a dedicated CSV artifact with per-sample runtime,
token counts, and quality-relevant outputs.
